\documentclass[UTF8,11pt]{ctexart}
\usepackage[a4paper,margin=24mm]{geometry}
\usepackage{amsmath,amssymb,amsthm,mathtools,mathrsfs}
\usepackage[all]{xy}
\usepackage{xurl,hyperref,enumitem}
\usepackage{tikz}
\usetikzlibrary{arrows.meta,cd,decorations.markings,calc,patterns}
\hypersetup{hidelinks,breaklinks=true}
\setlength{\emergencystretch}{3em}
\allowdisplaybreaks
\newtheorem*{theorem}{Theorem}
\newtheorem*{thm}{Theorem}
\newtheorem*{lemma}{Lemma}
\newtheorem*{proposition}{Proposition}
\newtheorem*{prop}{Proposition}
\newtheorem*{corollary}{Corollary}
\newtheorem*{cor}{Corollary}
\newtheorem*{claim}{Claim}
\theoremstyle{definition}
\newtheorem*{definition}{Definition}
\newtheorem*{defn}{Definition}
\newtheorem*{remark}{Remark}
\newtheorem*{example}{Example}
\newcommand{\R}{\mathbb R}
\newcommand{\C}{\mathbb C}
\newcommand{\Z}{\mathbb Z}
\newcommand{\Q}{\mathbb Q}
\newcommand{\N}{\mathbb N}
\newcommand{\F}{\mathbb F}
\newcommand{\D}{\mathbb D}
\newcommand{\bB}{\mathbb B}
\newcommand{\bC}{\mathbb C}
\newcommand{\bR}{\mathbb R}
\newcommand{\bZ}{\mathbb Z}
\newcommand{\bN}{\mathbb N}
\newcommand{\bQ}{\mathbb Q}
\newcommand{\bD}{\mathbb D}
\newcommand{\bF}{\mathbb F}
\newcommand{\CP}{\mathbb{CP}}
\newcommand{\RP}{\mathbb{RP}}
\newcommand{\sO}{\mathscr O}
\newcommand{\sI}{\mathscr I}
\newcommand{\sF}{\mathscr F}
\newcommand{\sG}{\mathscr G}
\newcommand{\sH}{\mathscr H}
\newcommand{\sR}{\mathscr R}
\newcommand{\sM}{\mathscr M}
\newcommand{\sV}{\mathscr V}
\newcommand{\sU}{\mathscr U}
\DeclareMathOperator{\Hom}{Hom}
\DeclareMathOperator{\Ext}{Ext}
\DeclareMathOperator{\Tor}{Tor}
\DeclareMathOperator{\Spec}{Spec}
\DeclareMathOperator{\Proj}{Proj}
\DeclareMathOperator{\supp}{supp}
\DeclareMathOperator{\im}{im}
\DeclareMathOperator{\id}{id}
\DeclareMathOperator{\Tr}{Tr}
\DeclareMathOperator{\tr}{tr}
\DeclareMathOperator{\rank}{rank}
\DeclareMathOperator{\ord}{ord}
\DeclareMathOperator{\dist}{dist}
\DeclareMathOperator{\End}{End}
\DeclareMathOperator{\Aut}{Aut}
\DeclareMathOperator{\Gal}{Gal}
\DeclareMathOperator{\vol}{vol}
\DeclareMathOperator{\Cl}{Cl}
\DeclareMathOperator{\coker}{coker}
\DeclareMathOperator{\codim}{codim}
\DeclareMathOperator{\divergence}{div}
\DeclareMathOperator{\Ric}{Ric}
\DeclareMathOperator{\grad}{grad}
\DeclareMathOperator{\Hess}{Hess}
\newcommand{\abs}[1]{\left\lvert#1\right\rvert}
\newcommand{\sabs}[1]{\lvert#1\rvert}
\newcommand{\babs}[1]{\bigl\lvert#1\bigr\rvert}
\newcommand{\Babs}[1]{\Bigl\lvert#1\Bigr\rvert}
\newcommand{\bbabs}[1]{\biggl\lvert#1\biggr\rvert}
\newcommand{\BBabs}[1]{\Biggl\lvert#1\Biggr\rvert}
\newcommand{\norm}[1]{\left\lVert#1\right\rVert}
\newcommand{\snorm}[1]{\lVert#1\rVert}
\newcommand{\bnorm}[1]{\bigl\lVert#1\bigr\rVert}
\newcommand{\Bnorm}[1]{\Bigl\lVert#1\Bigr\rVert}
\newcommand{\bbnorm}[1]{\biggl\lVert#1\biggr\rVert}
\newcommand{\BBnorm}[1]{\Biggl\lVert#1\Biggr\rVert}
\newcommand{\widebar}[1]{\overline{#1}}
\newcommand{\nicefrac}[2]{\frac{#1}{#2}}
\newcommand{\linnprod}[2]{\langle#1,#2\rangle}
\newcommand{\myindex}[1]{#1}
\newcommand{\glsadd}[1]{}
\newcommand{\avoidbreak}{}
\newcommand{\sourceref}[1]{\textnormal{[原文：\nolinkurl{#1}]}}
\newcommand{\thmref}[1]{\sourceref{#1}}
\newcommand{\lemmaref}[1]{\sourceref{#1}}
\newcommand{\propref}[1]{\sourceref{#1}}
\newcommand{\corref}[1]{\sourceref{#1}}
\newcommand{\defnref}[1]{\sourceref{#1}}
\newcommand{\exampleref}[1]{\sourceref{#1}}
\newcommand{\exerciseref}[1]{\sourceref{#1}}
\newcommand{\figureref}[1]{\sourceref{#1}}
\newcommand{\sectionref}[1]{\sourceref{#1}}
\newcommand{\appendixref}[1]{\sourceref{#1}}
\newcommand{\stacksref}[2]{\href{https://stacks.math.columbia.edu/tag/#1}{#2 [#1]}}

\begin{document}
\section*{047\quad 算术级数中素数的Dirichlet密度与L函数非消失}
\subsection*{来源与计数目标}
Andrew V. Sutherland，MIT 18.785 \emph{Number Theory I}，Fall 2021，Lecture 18。Proposition 18.20（p.~6）、\S18.4 完整推导与 Definition 18.22（pp.~7--8）。采用 MIT OpenCourseWare 实际访问版本；获取日期：2026-09-28。
\par\url{https://ocw.mit.edu/courses/18-785-number-theory-i-fall-2021/mit18_785f21_lec18.pdf}
Andrew V. Sutherland，MIT 18.785 \emph{Number Theory I}，Fall 2021，Lecture 19。\S19.3--19.4；Proposition 19.14、Theorems 19.15--19.16 及完整证明（pp.~10--12）。采用 MIT OpenCourseWare 实际访问版本；获取日期：2026-09-28。
\par\url{https://ocw.mit.edu/courses/18-785-number-theory-i-fall-2021/mit18_785f21_lec19.pdf}
\par\textbf{本文件只计一个问题：}证明互素剩余类含无穷多素数且 Dirichlet 密度为 $1/\varphi(m)$；核心非主 $L(1,\chi)\ne0$ 及 cyclotomic Euler 因子分解不能省去。
\subsection*{作者命题与人类证明}
\paragraph{Necessary notation (from Lectures 18--19).}
A Dirichlet character of modulus $m$ is the periodic extension to $\Z$ of a character of $(\Z/m\Z)^\times$, with value zero on integers not coprime to $m$. Write
\[
L(s,\chi)=\sum_{n\geq1}\chi(n)n^{-s}=\prod_p(1-\chi(p)p^{-s})^{-1}\qquad(\operatorname{Re}s>1).
\]
The set $X(m)$ consists of primitive Dirichlet characters of conductor dividing $m$; it is identified with the character group of $(\Z/m\Z)^\times$. Its product is obtained by taking the unique primitive character inducing the pointwise product. We have an isomorphism
\[
\varphi:\Gal(\Q(\zeta_m)/\Q)\xrightarrow{\sim}(\Z/m\Z)^\times,\qquad \sigma(\zeta_m)=\zeta_m^{\varphi(\sigma)}.
\]
For $H\subseteq X(m)$ put
\[
\phi(H):=\{\sigma\in\Gal(\Q(\zeta_m)/\Q):\chi(\sigma)=1\text{ for all }\chi\in H\},
\quad K=\Q(\zeta_m)^{\phi(H)}.
\]
Then $H$ is the character group of $\Gal(K/\Q)$ (Proposition 18.40).
\begin{proposition}[18.20]
Let $\chi$ be a non-principal Dirichlet character of modulus $m$. Then $L(s,\chi)$ extends to a holomorphic function on $\operatorname{Re}s>0$.
\end{proposition}
\begin{proof}
Define the function $T:\R_{\geq0}\to\C$ by $T(x):=\sum_{0<n\leq x}\chi(n)$. For any $x\in\R_{\geq0}$ Lemma 18.15 implies
\[
T(x+m)-T(x)=\sum_{x<n\leq x+m}\chi(n)=\sum_{n\in\Z/m\Z}\chi(n)=0,
\]
since $\chi$ is non-principal. Thus $T(x)$ is periodic modulo $m$ and therefore bounded.

Writing $L(s,\chi)$ as a Stieltjes integral (see \S18.5) and integrating by parts yields
\begin{align*}
L(s,\chi)&=\sum_{n\geq1}\chi(n)n^{-s}=\int_0^\infty x^{-s}\,dT(x)\\
&=\left.x^{-s}T(x)\right|_0^\infty-\int_0^\infty T(x)\,d(x^{-s})\\
&=0-\int_0^\infty T(x)(-sx^{-s-1})\,dx
=s\int_0^\infty T(x)x^{-s-1}\,dx.
\end{align*}
The RHS is holomorphic on $\operatorname{Re}s>0$, since it is the limit of the uniformly converging sequence of functions $\varphi_n(s):=s\int_0^nT(x)x^{-s-1}\,dx$ (here we use the fact that $T(x)$ is bounded), and is thus the analytic continuation of $L(x,\chi)$ to $\operatorname{Re}(s)>0$.
\end{proof}
\paragraph{19.3 Cyclotomic zeta functions and Dirichlet $L$-functions}
\begin{proposition}[19.14]
Let $p$ be a prime, let $m$ be a positive integer, and let $m'=m/p^{v_p(m)}$. Then $\Q(\zeta_{m'})$ is the maximal extension of $\Q$ in $\Q(\zeta_m)$ unramified at $p$. In particular, if $p$ does not divide $m$ then $\Q(\zeta_m)$ is unramified at $p$.
\end{proposition}
\begin{proof}
By Corollary 10.18, the extension $\Q_p(\zeta_{m'})/\Q_p$ is unramified. It follows from Proposition 12.4 that $\Q(\zeta_{m'})/\Q$ is unramified at $p$. Applying the same argument to all primes $q\ne p$ dividing $m$ shows that the extension $\Q(\zeta_{p^{v_p(m)}})$ is ramified only at $p$. By Corollary 14.27, there are no nontrivial unramified extensions of $\Q$, so every subfield of $\Q(\zeta_{p^{v_p(m)}})$ that properly contains $\Q$ is ramified at $p$. Now $\Q(\zeta_m)$ is the compositum of $\Q(\zeta_{p^{v_p(m)}})$ and $\Q(\zeta_{m'})$, which intersect in $\Q$, so any nontrivial extension of $\Q(\zeta_{m'})$ in $\Q(\zeta_m)$ contains a subfield of $\Q(\zeta_{p^{v_p(m)}})$ properly containing $\Q$ which must be ramified at $p$; the proposition follows.
\end{proof}
\begin{theorem}[19.15]
Let $H\subseteq X(m)$ be a group of primitive Dirichlet characters and let $K=\Q(\zeta_m)^{\phi(H)}$ be the corresponding subfield of $\Q(\zeta_m)$, with $\phi(H)$ defined as above. Then
\[
\zeta_K(s)=\prod_{\chi\in H}L(s,\chi).
\]
\end{theorem}
\begin{proof}
On the LHS we have
\[
\zeta_K(s)=\prod_{\mathfrak p}(1-N(\mathfrak p)^{-s})^{-1}
=\prod_p\prod_{\mathfrak p\mid p}(1-N(\mathfrak p)^{-s})^{-1},
\]
and on the RHS we have
\[
\prod_{\chi\in H}L(s,\chi)=\prod_{\chi\in H}\prod_p(1-\chi(p)p^{-s})^{-1}
=\prod_p\prod_{\chi\in H}(1-\chi(p)p^{-s})^{-1}.
\]
It thus suffices to prove
\[
\prod_{\mathfrak p\mid p}(1-N(\mathfrak p)^{-s})\stackrel{?}{=}\prod_{\chi\in H}(1-\chi(p)p^{-s})\tag{5}
\]
for each prime $p$.

Since $K/\Q$ is Galois, we have $[K:\Q]=e_pf_pg_p$, where $e_p$ is the ramification index, $f_p$ is the residue field degree, and $g_p=\#\{\mathfrak p\mid p\}$. On the LHS of (5) we have
\[
\prod_{\mathfrak p\mid p}(1-N(\mathfrak p)^{-s})=(1-(p^{f_p})^{-s})^{g_p}=(1-(p^{-s})^{f_p})^{g_p},
\]
which we note does not change if we replace $K$ with the maximal subfield $K'$ of $K$ in which $p$ is unramified (since $K/K'$ is totally ramified at every prime of $K'$ above $p$, only $e_p$ changes, not $f_p$ or $g_p$). On the RHS of (5), we have $\chi(p)=0$ for all $\chi\in H$ with conductor divisible by $p$, so we can replace $H$ with the subgroup $H'$ of Dirichlet characters with conductors prime to $p$. It follows from Proposition 19.14 that $K'=\Q(\zeta_m)^{\phi(H')}$ (to see this, note that if we put $m'=m/p^{v_p(m)}$ then $K'=K\cap\Q(\zeta_{m'})$ and $H'=H\cap X(m')$). Thus without loss of generality we assume $p\nmid m$, so $K$ is unramified at $p$ and we have $\#H=[K:\Q]=f_pg_p$.

Since $K/\Q$ is abelian and unramified at $p$, the Artin map gives us a Frobenius element $\sigma_p$ corresponding to the Frobenius automorphism $x\mapsto x^p$ of the residue field, which by definition has order $f_p$, so $\sigma_p$ has order $f_p$ in $\Gal(K/\Q)$. Viewing $H$ as the character group of $\Gal(K/\Q)$, the map $\chi\mapsto\chi(\sigma_p)$ defines a surjective homomorphism from $H$ to the group of $f_p$-th roots of unity $\alpha\in\mathrm U(1)$, and the kernel of this map has cardinality $\#H/f_p=g_p$. Therefore
\[
\prod_{\chi\in H}(1-\chi(p)p^{-s})=\prod_{\alpha^{f_p}=1}(1-\alpha p^{-s})^{g_p}=(1-(p^{-s})^{f_p})^{g_p},
\]
where the second equality follows from the identity $\prod_{\alpha^{f_p}=1}(1-\alpha T)=1-T^{f_p}\in\C[T]$.
\end{proof}
\paragraph{19.4 Non-vanishing of Dirichlet $L$-functions with non-principal character}
\begin{theorem}[19.16]
Let $\psi$ be any non-principal Dirichlet character. Then $L(1,\psi)\ne0$.
\end{theorem}
\begin{proof}
Let $\psi$ be a non-principal Dirichlet character, say of modulus $m$. Then $\psi$ is induced by a non-trivial primitive Dirichlet character $\widetilde\psi$ of conductor $\widetilde m$ dividing $m$. The $L$-functions of $\psi$ and $\widetilde\psi$ differ at only finitely many Euler factors $(1-\psi(p)p^{-s})^{-1}$ (corresponding to primes $p$ dividing $m/\widetilde m$), and these factors are clearly nonzero at $s=1$, since $p>1$. We thus assume without loss of generality that $\psi=\widetilde\psi$ is primitive.

Let $K$ be the $m$th cyclotomic field $\Q(\zeta_m)$. By Theorem 19.15 we have $\zeta_K(s)=\prod_\chi L(s,\chi)$, where $\chi$ ranges over the primitive Dirichlet characters of conductor dividing $m$, including $\psi$. By the analytic class number formula (Theorem 19.12), the LHS has a simple pole at $s=1$, and the same must be true of the RHS. Thus
\begin{align*}
\operatorname{ord}_{s=1}\zeta_K(s)&=\operatorname{ord}_{s=1}\prod_\chi L(s,\chi),\\
-1&=\operatorname{ord}_{s=1}L(s,1)\prod_{\chi\ne1}L(s,\chi),\\
-1&=\operatorname{ord}_{s=1}\zeta(s)\prod_{\chi\ne1}L(s,\chi),\\
-1&=-1+\sum_{\chi\ne1}\operatorname{ord}_{s=1}L(s,\chi).
\end{align*}
Each $\chi\ne1$ in the sum is necessarily non-principal (since it is primitive). We proved in Proposition 18.20 that for non-principal $\chi$ the Dirichlet $L$-series $L(s,\chi)$ is holomorphic on $\operatorname{Re}(s)>0$, thus $\operatorname{ord}_{s=1}L(s,\chi)\geq0$ for all $\chi$ appearing in the sum, which can therefore be zero if and only if every term $\operatorname{ord}_{s=1}L(s,\chi)$ is zero. So $L(1,\chi)\ne0$ for every non-trivial primitive Dirichlet character $\chi$ of conductor dividing $m$, including $\psi$.
\end{proof}
\paragraph{18.4 Primes in arithmetic progressions}
\begin{proof}[Author's proof in Section 18.4]
We now return to our goal of proving Dirichlet's theorem on primes in arithmetic progressions. It suffices to show that for any $a\perp m$ the sum $\sum_{p\equiv a\bmod m}p^{-s}$ is unbounded as $s\to1^+$. To convert this to a sum over all primes we use Proposition 18.37 to construct the indicator function
\[
\frac1{\varphi(m)}\sum_{\chi\in X(m)}\chi(p/a)=
\begin{cases}1&\text{if }p\equiv a\bmod m,\\0&\text{otherwise},\end{cases}
\]
where $p/a$ is computed modulo $m$ and $\chi$ ranges over primitive Dirichlet characters of conductor dividing $m$, which we identify with the character group of $(\Z/m\Z)^\times$ via Corollary 18.16. As $s\to1^+$ we have
\begin{align*}
\sum_{p\equiv a\bmod m}p^{-s}
&=\sum_p p^{-s}\frac1{\varphi(m)}\sum_{\chi\in X(m)}\chi(p/a)\\
&=\sum_{\chi\in X(m)}\frac{\chi(1/a)}{\varphi(m)}\sum_p\chi(p)p^{-s}\\
&=\sum_{\chi\in X(m)}\frac{\chi(1/a)}{\varphi(m)}(\log L(s,\chi)+O(1))\\
&=\frac{\log\zeta(s)}{\varphi(m)}+\sum_{\substack{\chi\in X(m)\\\chi\ne1}}\frac{\chi(1/a)}{\varphi(m)}\log L(s,\chi)+O(1).
\end{align*}
We now make the key claim that so long as $\chi$ is not principal, we have $L(1,\chi)\ne0$. This implies that $\log L(s,\chi)=O(1)$ as $s\to1^+$ and therefore
\[
\sum_{p\equiv a\bmod m}p^{-s}=\frac{\log\zeta(s)}{\varphi(m)}+O(1)
\]
is unbounded as $s\to1^+$, since $\zeta(s)$ is. Moreover, Mertens' second theorem implies
\[
\sum_{\substack{p\leq x\\p\equiv a\bmod m}}\frac1p\sim\frac{\log\log x}{\varphi(m)}.
\]
which proves that there are infinitely many primes $p\equiv a\bmod m$. We can makes this statement more precise by using the notion of Dirichlet density.
\end{proof}
\begin{definition}[18.22]
The Dirichlet density of a set of primes $S$ is given by
\[
d(S):=\lim_{s\to1^+}\frac{\sum_{p\in S}p^{-s}}{\sum_pp^{-s}},
\]
defined whenever this limit exists.
\end{definition}
The Dirichlet density of the set $S=\{p\equiv a\bmod m\}$ is
\[
d(S)=\lim_{s\to1^+}\frac{\sum_{p\in S}p^{-s}}{\sum_pp^{-s}}
=\lim_{s\to1^+}\frac{\log\zeta(s)/\varphi(m)}{\log\zeta(s)}=\frac1{\varphi(m)}.
\]
\subsection*{整理说明（非作者正文）}
开头记号据两讲定义摘编，非作者连续原段。先列18.20，再收19.14--19.16，最后返回18.4；这是显式的跨讲支撑材料合并，各证明内部顺序保留。18.4 中 ``key claim'' 对应随后一讲19.16，已经收入，未将未证 claim 当完成。

标准先修为有限阿贝尔群特征标正交性与对偶对应、cyclotomic Galois 群、局部不分歧分圆扩张、Frobenius 及 $efg$ 关系、$\Q$ 无非平凡处处不分歧扩张、解析类数公式（本库046）、Stieltjes 积分与 Euler 积对数高次项有界估计。主证明中取非主 $L$ 函数的非消失并未只作为先修；它的 Euler 因子比较、零极点阶数论证均在本文。原18.4的 primitive-character 指示函数写法及有限坏素数措辞照录；原19.16的 Euler 因子中写 $\psi$ 而非 $\widetilde\psi$ 亦照录。更强的实变量加权素数渐近原文用一句 Mertens 推论说明，本题只计已有展开论证的 Dirichlet 密度与无穷性，不以那句推论扩张计数目标。

已查看 Lecture18 pp.~7--8 与 Lecture19 pp.~10--11 原图；18的p.~6、19的p.~12图像接口失败，已取得其解析正文。难度依据是余类筛选、分圆域素理想分解和解析非消失的组合，不是套用现成 Dirichlet 定理。
\subsection*{可修改的简短标注（非作者正文）}
$c$：算术级数中素数的分布

$a$：特征标正交筛选；Euler 积；分圆域；Frobenius；分歧与剩余次数；Dedekind zeta 分解；解析类数公式；零极点阶数；Dirichlet 密度

\texttt{cross\_theory\_status = yes}。余类条件编码为有限群特征标和；分圆域的素理想分解决定 Euler 因子，从 Dedekind zeta 的简单极点推出各非主 L 函数非消失，最终返回素数余类密度。位置：19.15--19.16 与18.4。

全部标注均为非穷尽、可修改初稿；未标注不表示未使用。
\subsection*{许可与修改记录}
材料来自 MIT OpenCourseWare，作者与课程见来源。按该站所示 CC BY-NC-SA 4.0 许可复用，整理部分沿用同一许可。\url{https://creativecommons.org/licenses/by-nc-sa/4.0/}。2026-09-28：助手转录题证、调整数学排版并加入来源、范围和初稿标注；不暗示作者或 MIT 认可本次整理。所留为本次题证转录摘录，不是原 PDF 下载件。
\end{document}
